\section{Data and Experimental Design}
\label{sec:data}

\subsection{Asset universe and calendar}

The primary file contains 2,946 daily price observations for seven proxies:
the S\&P 500, NASDAQ, Dow Jones Industrial Average, SSE 50, a China dividend
index, ChiNext, and gold. This produces 2,945 validated daily log returns.
A publication appendix must still document exact tickers, provider, currency,
time-zone close, price-adjustment convention, retrieval date, and the treatment
of non-synchronous US and Chinese closes.

A calendar helper constructs pairs consisting of a month-end decision date and
its next available month-end holding-period end. The last price observation is
therefore never treated as a decision. After the 250-day minimum history, there
are 114 training holding periods and exactly 24 locked evaluation periods. The
first evaluation decision is 31 July 2024 and its return ends 30 August 2024;
the final decision is 30 June 2026 and its shortened holding period ends
6 July 2026. This partial final period must be reported explicitly and checked
in a complete-month robustness sample when more data become available.

\subsection{Locked split and baseline configuration}

No evaluation observation may affect marginal fitting, empirical transforms,
vine order or parameters, neural early stopping, state normalization, RL
training, hyperparameter choice, benchmark choice, or diagnostic thresholds.

\begin{table}[H]
\centering
\caption{Locked baseline configuration}
\label{tab:configuration}
\begin{tabular}{ll}
\toprule
Component & Setting \\
\midrule
Assets & 7; net exposure 1.0, gross cap 1.5 \\
Rebalancing and decision horizon & Monthly; 24 actions \\
Recurrent context & 30 strictly prior monthly states \\
Locked holdout & Final 24 non-empty holding periods \\
Marginals & Diagnostic-gated GARCH grid; component-EWMA fallback \\
Vine & Training-ordered, six-tree Student-$t$ D-vine \\
Dynamic dependence & Neural correlation on all 21 edges \\
Synthetic pre-training & 1,000 paths, empirical monthly marginals, AR copula \\
Historical fine-tuning & 61 trajectories, eight balanced passes (488 episodes) \\
CVaR scenarios per state & 512; robustness at 1,024 and 2,048 \\
RL optimisation & Recurrent TD3; 64 hidden units, one LSTM layer \\
Costs & Turnover 0.001; annual short 0.03; cash borrowing 0.02 \\
\bottomrule
\end{tabular}
\end{table}

\subsection{Synthetic-data acceptance}

Synthetic data are a model diagnostic, not a performance result. Diagnostics
are calculated on monthly log returns. Marginal checks use standardised errors
for means and 5\% tail summaries rather than relative errors in gross means near
one. Pairwise correlation compatibility uses a moving-block bootstrap from the
historical training months. Five-percent lower-tail co-exceedance reports the
number of conditioning and joint events and exact binomial intervals; an
observed zero is not interpreted as asymptotic tail independence. Temporal
checks compare lag-one returns and squared returns. The generator writes all
diagnostics and exits non-zero if any required gate fails.

\subsection{Comparison and inference design}

The final evaluator accepts strategy weights, not self-reported wealth. It
applies the same realised assets, net/gross constraints, turnover, short fees,
and financing to every method. The minimum benchmark set is pre-declared:
equal weight, constrained sample mean--variance, DCC-GARCH, static vine,
rolling vine, no-vine TD3, no-pre-training TD3, and the risk-neutral ablation.
Legacy benchmark artifacts are excluded.

Development uses expanding-origin pre-holdout folds and validation windows.
The final 24 periods are inspected once after configurations are frozen.
Report all training seeds, compute time, CAGR, arithmetic excess-return Sharpe,
volatility, maximum drawdown, 5\% CVaR, terminal wealth, turnover, and short
notional. Paired HAC utility tests with Holm correction and a moving-block
Reality Check address serial dependence and data snooping. Twenty-four months
nevertheless have low power; effect sizes and intervals are primary.

